%root = ../thesis-main.tex
\chapter{Reactivity and Dependency Management}\label{ch:reactvity_dep_mgmnt}

The previous chapter introduced what is a simulator, with an emphasis to
the discrete-event formalism. Moreover, the key concept of dependency management
is introduced in such systems with emphasis on complex systems with sparse
dependencies such as biochemical systems kinetic. This chapter aims to introduce
a complementary viewpoint in dependency management and resolution through
\emph{reactivity} as a general-purpose technique for expressing and executing
such dependency structures.

Reactive programming is often present as a software-engineering style for
building responsive user interfaces or streaming
applications~\cite{kuhn2017RDP}, but at more abstract level, however, it
provides a vocabulary and a set of execution strategies and principles for
maintaining \emph{derived values} and \emph{event-driven processes} consistent
with a changing base state (e.g. a simulated environment that changes over
time). In other words, it can be interpreted as a dependency resolution
mechanism: changes propagate through a set of dependencies in an implicit,
dynamic graph triggering recomputation, notifications and potentially the
scheduling of future work.

In what follows, an extensive theoretical foundation is built around three
complementary formalisms:
\begin{itemize}
	\item \emph{event-riven} and \emph{data-flow} models, which view
	      computation as the propagation of tokens or events through a network of
	      operations;
	\item the \emph{observer} family of patterns, traditionally introduced in
	      object-oriented design, which popularised callback-based push
	      propagation and decoupled notification;
	\item \emph{Functional Reactive Programming}, which provides specific
	      meanings and manipulation tools to time-varying values and event
	      streams.
\end{itemize}
We then connect these ideas to modern reactive frameworks (with Project Reactor
as a concrete reference point) and highlight the implications for designing
reactive dependency management in a discrete-event simulation engine.

\section{Theory of Reactive Systems}

\subsection{Event-driven and Data-flow Computation}

A natural conceptual root of reactive programming is \emph{event-driven}
computation. In such systems, execution follows a data-driven model, where each
instruction becomes executable only once all required operands have been
preceding instructions~\cite{dennis1974PABDFP}. Applied to software systems,
such formalism define a system organised as a graph of components whose
dependencies are realised through the flow of messages. Systems developed under
such assumptions are usually employed to exploit parallelism and concurrency,
especially for distributed systems, which introduces a higher degree of
complexity due to how communication is handled among the nodes of the system.
Typical event-driven applications rely on \emph{asynchronous communication},
which avoids keeping nodes in hold when sending messages while using explicit
mechanism of synchronisation when needed~\cite{lamport1978TCOEDS}. Conversely,
\emph{synchronous communication} blocks whenever a message is sent. No approach
excludes the other, and they could be better suited depending to the problem at
hand (see \cite{arjomandi1983ESVADS} for a case where synchronous systems could
have better performance). A perfect exemplification of such architecture at the
programming-language level, is the \emph{actor model} which is mainly based on
message-passing. In the actor model each actor encapsulates state and
behaviour, processes one message at a time, and interacts exclusively by
sending messages to other actors. This yields a compositional model where
concurrency emerges from independent actors \emph{reacting} to incoming
messages~\cite{agha1986AMCCDS}.

\subsection{Reactive systems}
The term \emph{reactive systems} was originally coined to distinguish systems
that  maintain a continuous interaction with their environment from
transformational systems that simply compute a result from
inputs~\cite{harel1989DRS}. In modern software engineering however, this notion
denotes systems that must continuously respond to incoming stimuli (user
actions, network messages, sensor readings) under certain constraints of
latency and availability. The \emph{Reactive Manifesto}\footnote{
	\url{https://reactivemanifesto.org/}} organises this space under four
high-level properties: responsiveness, resilience, elasticity and
message-driven interaction.

Although the manifesto is not a formal model of computation, it highlights
constraints that directly affect dependency resolution mechanisms:
\begin{itemize}
	\item \emph{backpressure} and load regulation which determine how fast
	      communications along dependencies edges are allowed to propagate;
	\item \emph{asynchrony} and distribution to maximise responsiveness,
	      but making the arrival of updates uncertain and complicating consistency;
	\item \emph{failure} as a first-class concern, which motivates explicit lifecycle
	      management for subscription and resources.
\end{itemize}
These considerations will become relevant when discussing about propagation strategies
and pitfalls of reactive frameworks. [LINKS TO SECTION HERE AND ABOVE]

\subsection{The Observer pattern and push-based propagation}
In object-oriented systems, reactivity is traditionally implemented through the
\emph{Observer pattern}~\cite{gamma1994design}. A \emph{subject} maintains a
set of \emph{observers} and notifies them whenever its state changes.
Therefore, the pattern decouples producers from consumers: the subject does not
need to know what observers do with the updates, only that they implement a
common notification interface.

From a dependency-management viewpoint, observer-style reactivity correspond to
a \emph{push-based} propagation, where updates originate at a source which
\emph{immediately} (\emph{synchronous}) push them to dependents which may in
turn update their own status.

This approach is simple and effective, but it introduces well-known challenges
to the point that some may argue to ``Deprecate the Observer pattern''~\cite{maier2010DOP}.
Among some notable issues we find
\begin{enumerate*}[label=(\roman*)]
	\item the dependency graph is often implicit in registrations spread across
	      the codebase
	\item ordering constraints are subtle when several observers depend on one
	      another
	\item feedback cycles and re-entrant notifications can create inconsistent
	      intermediate state (link to glitches) and
	\item lifecycle management becomes critical (link to section)
\end{enumerate*}

These issues motivate richer reactive abstractions where dependencies become
explicit and where the runtime can provide stronger guarantees about lifecycles
and propagation order.

\subsection{Functional Reactive Programming}
\ac{FRP} emerged as an attempt to provide a more principled account of
reactive computation. In the classical \ac{FRP} formulation, time-varying
values (often referred to as \emph{behaviours} or \emph{signals}) are treated
as first-class values and discrete occurrences are represented as \emph{event
streams}~\cite{elliott1997FRA}. 

